\specialsection{Заключение}

В дипломной работе была рассмотрена задача рекомендации следующей культуры для конкретного сельскохозяйственного поля на
основе его исторической последовательности использования и базовых признаков поля и региона. Работа изначально строилась
не как попытка создать полный оптимизатор севооборота, а как исследовательский прототип объяснимой рекомендательной
системы, ориентированной на поддержку принятия решений. Такой выбор постановки позволил сосредоточиться на построении
воспроизводимого контура анализа, в котором сочетаются предиктивное ядро, рекомендательный режим, оценка уверенности и
интерпретационный слой.

В ходе работы были подготовлены и структурированы данные USDA NASS CSB/CDL, выполнено сопоставление и укрупнение
культурных классов, построена оконная выборка вида \texttt{history\_1,\ history\_2,\ history\_3\ -\textgreater{}\
target}, а также зафиксировано воспроизводимое разделение данных на обучающую, валидационную и тестовую выборки. Это
позволило сформировать единую основу для всех последующих этапов исследования и обеспечить корректное сравнение
различных подходов.

На этапе моделирования были рассмотрены простые baseline-стратегии, вероятностный ориентир Markov-1 и модель CatBoost.
Сравнение показало, что финальная baseline-модель CatBoost заметно превосходит более простые подходы по основным
классификационным и рекомендательным метрикам. Это подтвердило, что использование трёхлетней истории и контекстных
признаков даёт существенный прирост качества по сравнению с моделями, опирающимися только на частотные или одношаговые
переходные закономерности.

Одним из ключевых результатов работы стал переход от обычной многоклассовой классификации к рекомендательной постановке.
Было показано, что наиболее естественным прикладным форматом использования модели является не top-1 предсказание, а
Top-k shortlist вероятных следующих культур. Полученные результаты продемонстрировали, что даже в тех случаях, когда
первая рекомендация оказывается недостаточно точной, правильный класс часто присутствует в числе нескольких наиболее
вероятных кандидатов. Это позволяет рассматривать систему именно как рекомендательный инструмент, а не как жёсткий
одновариантный классификатор.

Важной частью работы стал confidence-анализ. Введение confidence-слоя позволило показать, что рекомендации модели
неоднородны по степени надёжности и естественным образом разделяются на более устойчивые и более спорные случаи. Анализ
показал, что в группе высокой уверенности модель существенно точнее, чем в среднем по выборке, а в случаях низкой
уверенности сохраняет практическую полезность прежде всего как shortlist-рекомендатель. Дополнительный анализ калибровки
подтвердил, что вероятности модели могут использоваться не только для ранжирования кандидатов, но и как содержательный
confidence-сигнал.

Следующим значимым результатом стало построение explainability-слоя на основе knowledge graph и набора интерпретационных
правил. В работе показано, что knowledge graph целесообразно использовать не как вторую предиктивную модель и не как
полноценную экспертную систему, а как семантический слой для интерпретации shortlist-кандидатов. Такой подход позволил
дополнить модельный вывод support- и warning-сигналами, графовой поддержкой и кратким объяснительным контекстом.
Особенно важным это оказалось для сложных и неоднозначных случаев, в которых полезность shortlist проявляется сильнее
всего.

Для демонстрации практического сценария использования результатов был разработан пространственный демонстрационный слой.
Он связал рекомендации и explainability-контекст с конкретными полями на карте и показал, как модельный результат может
быть представлен в форме пользовательского интерфейса. При этом spatial demo не изменяет предиктивное ядро системы, а
выступает как аналитическая и интерфейсная надстройка, делающая результаты более наглядными и удобными для
интерпретации.

Таким образом, в работе был построен целостный воспроизводимый прототип рекомендательной системы по севообороту,
включающий:

\begin{itemize}
\tightlist
\item
  подготовку данных и построение временной выборки;
\item
  baseline-сравнение моделей;
\item
  финальное предиктивное ядро на основе CatBoost;
\item
  рекомендательный Top-k режим;
\item
  confidence-layer;
\item
  explainability-слой на основе knowledge graph;
\item
  пространственную демонстрацию результатов.
\end{itemize}

Основной итог работы состоит в том, что задача выбора следующей культуры может быть продуктивно рассмотрена в формате
explainable recommendation / decision-support prototype. В этой постановке система не подменяет экспертное решение, но
помогает пользователю анализировать вероятные варианты, видеть уровень уверенности модели и получать интерпретируемый
контекст для shortlist-кандидатов. Именно в таком виде разработанный прототип представляет практический и
исследовательский интерес.

В качестве направлений дальнейшего развития могут рассматриваться расширение признакового пространства за счёт внешних
агрономических факторов, углубление confidence-анализа, развитие knowledge graph и правил интерпретации, а также
дальнейшее развитие пространственного интерфейса. Однако даже в текущем виде работа демонстрирует завершённый и
методически согласованный контур решения, достаточный для того, чтобы рассматривать её как полноценный результат
дипломного исследования.
